\chapter{Limitations and Future Work}
\label{ch:limitations-future-work}

This chapter states what this thesis does not establish. It is written on the principle that a
limitation named by the author is worth more than one found by the reader, and it includes the
hypotheses that were registered before data collection and then never tested --- reported as
such rather than dropped, because a hypothesis that vanishes between registration and reporting
is indistinguishable, to a reader, from one that was tested and inconvenient.

\section{Limitations}
\label{sec:limitations}

\subsection{The Larger Corpus Is Synthetic and Born-Digital}
\label{sec:lim-corpus}

The 146-document corpus is generated: its pages are rendered from the same data that supplies
the reference, so they are clean by construction --- no skew, no shadow, no stamp, no staple, no
handwriting, no compression artefact. Real documents in a tax firm arrive with all of those.

The held-out corpus exists precisely to bound this, and it does so with a number:
phone-scanned documents score more than eleven points of mean per-invoice F$_1$ below
email-native ones (Table~\ref{tab:heldout-by-channel}). \textbf{Any figure from the synthetic
corpus should therefore be read as an upper bound for the camera-captured case}, and the size of
the correction is \emph{at least} that gap --- the scans here are deskewed scan-app output, the
mildest form of camera capture (\S\ref{sec:method-data}). The direction is corroborated
independently by a multilingual parsing benchmark reporting an average 17.8\,\% drop for
open-source parsers on photographed documents \cite{mdpbench2026}.

\subsection{The Held-Out Corpus Is Small and Single-Annotator}
\label{sec:lim-heldout}

Thirty-nine documents, adjudicated by one person. Two consequences follow.

\textbf{No precise figure is supported.} The headline held-out number should be read as
establishing a level and a direction, not a value to two decimal places. The per-channel cells
are smaller still, some containing around ten documents, so the channel comparison carries the
weight of a clear effect on a small sample rather than of a precise estimate.

\textbf{No inter-annotator agreement can be reported.} The three-channel design substantially
limits the exposure --- 463 accepted cells rest on printed-text proof or on two independent
channels agreeing, not on judgement (\S\ref{sec:method-heldout-gt}) --- but 248 cells were
decided by the author, and for those there is no second opinion. This is irreducible for a
single-author corpus and is the reason the datasheet in Appendix~\ref{app:heldout} records the
provenance class of every cell rather than merely asserting that the key was reviewed.

\subsection{Header Fields Only, on Real Documents}
\label{sec:lim-groups}

The real-invoice result covers flat header fields. Repeating groups --- line items, the tax
breakdown per rate, discount terms --- are excluded structurally, not by a reporting convention.

The reason is stated in \S\ref{sec:method-heldout-gt}: rows do not align across reading channels,
because one reader splits a line another merges, so positional comparison manufactures
disagreement out of segmentation rather than out of error. The group rows in the held-out key were
seeded from a single channel and never author-reviewed, and scoring against an unreviewed key is
what produced the retracted figure this project already had to withdraw once.

Line-item extraction is therefore measured on the synthetic corpus, where the reference is exact
and free (Table~\ref{tab:attribution-clusters}). The gap that remains is a measurement of
\emph{real} line items, and \S\ref{sec:fw-groups} states the route to it.

\subsection{The Structuring Model Was Held Constant, Not Selected}
\label{sec:lim-structurer-selection}

This is the sharpest asymmetry in the thesis and it deserves to be stated plainly.

The reading model was chosen competitively: a cohort survey, then a shortlist, then a direct
comparison on two instruments, then a hand-judged audit of all 52 residual misses, then a
pre-registered conjunctive test of the larger sibling that the challenger failed
(\S\ref{sec:results-reader}). \textbf{The structuring model was not chosen at all.} It was fixed
as the controlled variable so that the reader comparison and the specification-repair work could
be interpreted against a stable back end, and no structuring bake-off was ever run.

Holding it constant is what makes the rest interpretable --- swapping it would have invalidated
every committed baseline, and the attribution result depends on the structuring stage being the
same in both arms. But the limitation is real and it bounds two claims. No claim can be made that
this is the best available structuring model, and the adaptation results of
\S\ref{sec:results-adaptation} are conditional on it: a different structuring model might close
the reader-noise gap with no adaptation at all.

There is a related asymmetry worth reporting, because it is the reason the adaptation experiment
was designed as a grid rather than a before-and-after. On text rendered from the reference, this
structuring model already scores in the high nineties (Table~\ref{tab:sealed-val-arms}) --- above
the threshold the project had set. The measured shortfall is therefore \emph{robustness to
reading noise}, not capability, and an intervention aimed at capability was always unlikely to
help.

\subsection{One Adaptation Recipe, No Sweep}
\label{sec:lim-finetune}

The adaptation study tested one recipe: one rank, one learning rate, one schedule, one epoch
budget, 100 fitting examples (Table~\ref{tab:hyperparameters}). No hyperparameter search was
performed, and the training set contained no examples teaching appropriate abstention.

The negative result is therefore about \emph{this} recipe and not about adaptation as a
technique. \S\ref{sec:disc-finetune} sets out why the observed failure mode --- rising spurious
emission rather than degraded placement --- is consistent with published findings, and
\S\ref{sec:fw-abstention} states the specific untested variant that the mechanism points to.

\subsection{Hardware Envelope}
\label{sec:lim-hardware}

All inference was developed and measured inside a single laptop with 16\,\ac{GB} of unified
memory. That is a narrower claim than \emph{runs locally}: the privacy premise is on-premises
processing, which would also permit a workstation with a discrete accelerator, and such a machine
would likely support a larger reading model than the one selected here.

The envelope bound training rather than inference. The adaptation runs exceeded it and were moved
to a rented accelerator (\S\ref{sec:method-repro}), which means the local-only property holds for
the delivered pipeline but not for the process that produced its adapters. Since those adapters
are not delivered, the delivered system remains fully local --- but a firm wishing to adapt the
pipeline to its own documents on its own hardware could not do so within this envelope.

\subsection{Language, Document Class and Jurisdiction}
\label{sec:lim-scope}

German and English business-to-business invoices, under German and European invoicing
conventions. Nothing is claimed about receipts, contracts, bank statements, dunning letters or
correspondence; nothing about jurisdictions with different mandatory particulars; nothing about
languages outside the two measured. The single French-language document in the synthetic corpus is
an instance, not evidence of transfer.

\subsection{Measurement Constructs}
\label{sec:lim-measurement}

Two constructs carry more inferential weight than their definitions strictly support.

\textbf{The reading-quality proxy is a proxy.} It asks whether a reference value appears in the
transcript at all. A transcript containing every value on the page in scrambled order would score
perfectly and be useless to the structuring stage (\S\ref{sec:method-answerability}). The proxy
was consequently never used as the sole basis for a decision, and the one decision where it
mattered was settled by adjudicating individual misses.

\textbf{F$_1$ over header fields is not a measure of legal usability.} It measures extraction
correctness. It does not measure whether a booking would survive audit, whether the mandatory
particulars of \S14~\ac{UStG} are satisfied, or whether an input-tax deduction is defensible. The
weighted compliance metric that would have gestured toward this was abandoned deliberately
(\S\ref{sec:method-metrics}), and no legal assessment is offered anywhere in this thesis.

\subsection{Two Known Defects Remain Open}
\label{sec:lim-open-defects}

Reported here so that the apparatus is not represented as cleaner than it is.

\textbf{One field regressed under a correct repair and stays regressed.} Making the
perfect-text renderer faithful to the reference exposed a comparison problem in the supplier
article number field that was not resolved (\S\ref{sec:validity-open}).

\textbf{A stale configuration default remains a hazard.} A diagnostic that constructs the
experiment configuration from its defaults, rather than loading the project's configuration file,
silently selects a superseded reading model. The two audit tools were fixed to load the
configuration explicitly; the default itself was not changed, so any future script that builds
the configuration bare will measure the wrong pipeline (\S\ref{sec:validity-specification}).

\subsection{The Evaluation Methodology Is Not Itself Deployable}
\label{sec:lim-cloud-gt}

Building the held-out answer key required two cloud services, so the content of those 39 invoices
left local hardware (\S\ref{sec:method-heldout-gt}). The distinction relied upon --- a measuring
instrument is not part of the delivered system --- holds for the pipeline, which makes no network
call on any inference path. It does not hold for the methodology: a practitioner bound by
professional secrecy could not construct this answer key from their own client documents.
\S\ref{sec:fw-local-gt} states what closing that would require.

\section{Pre-Registered but Unevaluated Hypotheses}
\label{sec:unevaluated}

The scope of this thesis was narrowed during the work, in favour of one complete and honestly
measured layer over three partially measured ones. Six registered hypotheses were consequently
never tested. They are listed here as \emph{not evaluated}, which is distinct from refuted and
distinct from inconclusive: no data was collected against them, so they remain open.

\begin{enumerate}
  \item \textbf{Knowledge-graph fidelity.} That extracted records could be assembled into a
        graph whose entity resolution is accurate enough for cross-document reasoning. The layer
        was designed (\S\ref{sec:design-layer2}) and not built.
  \item \textbf{Graph versus flat retrieval.} That graph-structured retrieval would outperform
        flat vector retrieval on multi-document analytical questions. Untested: it presupposes
        the graph.
  \item \textbf{Query routing.} That a router could reliably dispatch a natural-language question
        to the appropriate retrieval strategy. Untested for the same reason.
  \item \textbf{Conditional validator retry.} That re-prompting the structuring stage with the
        specific validation failure it produced would recover a measurable share of failures. The
        validate-and-repair stage was built deterministic and single-pass; the retry variant was
        never implemented, and no figure in this thesis reflects one.
  \item \textbf{Cloud versus local comparison.} That the accuracy cost of the local-only
        constraint could be quantified against a frontier cloud model on the same corpus.
        Deliberately not run on client documents, for the reason the thesis exists;
        \S\ref{sec:fw-cloud} describes the form it could safely take.
  \item \textbf{Template shift.} That performance would degrade measurably on invoice layouts
        drawn from vendors absent from the fitting set. Partially and accidentally addressed by
        the held-out corpus, whose layouts were all unseen --- but not in the controlled form
        registered, since layout novelty and capture channel vary together there and cannot be
        separated.
\end{enumerate}

\section{Future Work}
\label{sec:future-work}

The items below are ordered by what the results of this thesis suggest, not by convenience. The
first two follow directly from the finding that reading is the binding constraint.

\subsection{Improve the Reading Stage, Not the Structuring Stage}
\label{sec:fw-reader}

The attribution result (\S\ref{sec:results-attribution}) says the largest available gain is in
reading. Three specific directions follow.

\textbf{Adapt the reading model rather than the structuring model.} This thesis adapted the stage
that was not the bottleneck. Low-rank adaptation of the vision tower on invoice pages is the
untested experiment that the attribution evidence actually recommends, and it is a different
experiment from the one reported here.

\textbf{Ensemble across readers with complementary failure modes.} The two finalists fail
differently and in a way that looks exploitable: one drops letterhead and footer regions while
reading characters near-flawlessly, the other reads every region with occasional character slips
(\S\ref{sec:results-reader}). A combination that takes regions from the first and characters from
the second is a concrete proposal, and the per-miss classification data needed to evaluate it
already exists.

\textbf{Pre-process before reading.} Deskewing, contrast normalisation and resolution
enhancement address the channel gap directly. Given that the channel gap is the dominant
real-world cost measured here, an intervention on the image may be worth more than any
intervention on the models.

\subsection{Measure Real Line Items}
\label{sec:fw-groups}

The route to closing \S\ref{sec:lim-groups} is known and is not research: adjudicate the held-out
group rows under a \emph{set-based} matching rule rather than a positional one, so that a reader
which merges two lines is not penalised as though it had misread them. That requires a matching
criterion defined before adjudication begins, and it requires the same ranked-escalation and
provenance-class discipline the header key used. It is expensive but it is bounded work.

\subsection{Teach Abstention Rather Than Always Answering}
\label{sec:fw-abstention}

The adaptation failed by raising spurious emission (\S\ref{sec:disc-finetune}), which is what
supervised fine-tuning on exclusively-filled targets would be expected to do
\cite{uluoglakci2026humility}. The corresponding untested variant is a training set that includes
appropriate absence as a taught behaviour: documents where a field is genuinely missing, with the
correct answer being the null. This thesis's training data contained no such examples, and the
mechanism points directly at that gap.

\subsection{Knowledge-Graph Construction}
\label{sec:fw-kg}

The most direct continuation of the designed-but-unbuilt architecture
(\S\ref{sec:design-layer2}). The interesting question is not whether a graph can be built --- it
can --- but how extraction error propagates into it: does a field-level error rate predict a
compliance-decision error rate, or does the graph's redundancy absorb some errors while amplifying
others? That is a measurement question, and the per-field outcome data this thesis produces is the
right input to it.

\subsection{Analytical Query and Retrieval Comparison}
\label{sec:fw-query}

Building the query layer (\S\ref{sec:design-layer3}) would let the second and third unevaluated
hypotheses be tested. The open question worth answering is narrower than the design: which
\emph{classes} of question genuinely benefit from graph structure over flat retrieval, since the
answer is unlikely to be uniform across question types.

\subsection{A Cloud Comparison That Respects the Premise}
\label{sec:fw-cloud}

Quantifying the accuracy cost of the local-only constraint requires running the same corpus
through a frontier cloud model. That cannot be done with client documents without violating the
premise. It can be done on the synthetic corpus, where every document is generated and none is
confidential --- and since the synthetic corpus carries exact reference data, such a comparison
would be cleanly scored. The result would be a lower bound on the cost of the constraint, because
the synthetic corpus is the easy case for both sides.

\subsection{Locally Adjudicable Ground Truth}
\label{sec:fw-local-gt}

Closing \S\ref{sec:lim-cloud-gt} requires the three-channel adjudication to run without a cloud
service: a local vision model as one channel, a local \ac{OCR}-plus-layout stack as a second, the
embedded text layer as a third, and the same printed-text gate and ranked escalation. Whether
three \emph{local} channels are independent enough for the two-channel-agreement rule to carry
its current warrant is itself an open question, and it is the one that would have to be answered
first.

\subsection{Robustness on Degraded Documents}
\label{sec:fw-robustness}

The channel gap measured here is the most actionable finding in the thesis, and it was measured on
ten photographed documents. A controlled degradation study --- the same documents captured at
several resolutions, angles and lighting conditions --- would turn a two-level comparison into a
response curve, and would say where on that curve a firm's document-handling practice needs to
sit. The protocol from the multilingual parsing benchmark \cite{mdpbench2026} is a reasonable
template.

\subsection{Production Hardening}
\label{sec:fw-production}

What a firm would additionally need: throughput and batching, audit logging of every extraction
with its provenance, and the review interface that the precision-over-recall error profile
(\S\ref{sec:results-heldout}) makes viable.

One specific and cheap addition belongs here. The validate-and-repair stage does not currently
check the arithmetic identities the invoice standard guarantees --- line amounts summing to the
line total, taxable base plus tax equalling the grand total, the amount due following from the
grand total and any prepayment (\S\ref{sec:design-validation}). Those identities are
deterministic, require no model, and would convert a class of silent misreading into a flagged
document. They would not improve any score reported in this thesis, which is precisely why they
were not built; they would improve the reviewability of the output, which is what a firm would
actually be buying. The prototype of
Chapter~\ref{ch:implementation} includes a review surface for exactly this reason, but it was
built to support the evaluation rather than a firm's daily workflow, and the difference between
those two is substantial.
